\lyseis{3.3}{Υπολογιστής τσέπης}
\sol{1} α) $365{,}852$ \quad β) $912{,}346$ \quad γ) $75{,}03$ \quad δ) $570$.
\sol{2} α) $50\cdot6=300$, άρα $287{,}33$. β) $20:5=4$, άρα $4{,}4$.
\sol{3} α) \key{(}\key{4.5}\kpl\key{2.7}\key{)}\kpo\key{1.25}\kis\ $\to9$. β) \key{36}\kdi\key{(}\key{0.9}\kpl\key{2.1}\key{)}\kis\ $\to12$. (Αν δεν έχει παρενθέσεις: πρώτα $0{,}9+2{,}1=3$ και μετά $36:3$.)
\sol{4} α) $6{,}5\cdot4-3{,}2\cdot5=26-16=10$. β) $(1{,}2+0{,}8):0{,}5=2:0{,}5=4$.
\sol{5} \key{MC} \key{2.5}\kpo\key{8}\kis\key{M+} \key{1.5}\kpo\key{6}\kis\key{M+} \key{4.25}\kpo\key{2}\kis\key{M$-$} \key{MR}. Αποτέλεσμα $20+9-8{,}5=20{,}5$.
\sol{6} α) $10-2{,}5=7{,}5$ \quad β) $12{,}6:2{,}1=6$ \quad γ) $0{,}25\cdot6{,}8+1{,}2=1{,}7+1{,}2=2{,}9$.
\sol{7} α) Λάθος: κάνει σωστά ό,τι του πατήσουμε, και το δεδομένο μπορεί να είναι λάθος. β) Λάθος: δείχνει προσέγγιση. γ) Σωστό. δ) Λάθος: τα απλά εκτελούν με τη σειρά πληκτρολόγησης. ε) Λάθος: η εκτίμηση βρίσκει τα λάθη πληκτρολόγησης.
\sol{8} $57{,}6-18{,}35=39{,}25$ €. Εκτίμηση: $58-18=40$.
\sol{9} Έβαλε $67{,}5$ αντί για $6{,}75$ (η υποδιαστολή έφυγε μία θέση δεξιά): $67{,}5\cdot1{,}2=81$. Το σωστό είναι $6{,}75\cdot1{,}2=8{,}1$. Εκτίμηση: $7\cdot1=7$.
\sol{10} $312:100\cdot7{,}8=3{,}12\cdot7{,}8=24{,}336$ λίτρα, δηλαδή $24{,}3$ λίτρα.
\sol{11} Η οθόνη δείχνει $14{,}285714$ (ή $14{,}285714286$ με περισσότερα ψηφία). Στρογγυλοποίηση: $14{,}29$ €. Συνολικά $7\cdot14{,}29=100{,}03$ €, δηλαδή τρία λεπτά παραπάνω από το ταμείο.
\sol{12} $0{,}06\cdot4{,}5=0{,}27$ kWh την ημέρα. Σε $30$ ημέρες $8{,}1$ kWh. Κόστος $8{,}1\cdot0{,}18=1{,}458\approx1{,}46$ €.
\sol{13} Το $1:3$ στην οθόνη δείχνει $0{,}33333333$ (προσέγγιση). Άρα $3\cdot0{,}33333333=0{,}99999999\neq1$. Η ακριβής τιμή $\frac{1}{3}$ έχει άπειρα ψηφία.
\sol{14} Εκτέλεσε τις πράξεις με τη σειρά που πατήθηκαν: $2+3=5$ και μετά $5\cdot4=20$. Για να βγει το σωστό πατάω πρώτα το γινόμενο: \key{3}\kpo\key{4}\kpl\key{2}\kis{} (δίνει $12+2=14$), ή υπολογίζω το $3\cdot4$ και το φυλάω με \key{M+}.
\sol{15} Ο $0{,}98$ είναι μικρότερος από το $1$, άρα το $487\cdot0{,}98$ είναι λίγο μικρότερο από το $487$, δηλαδή περίπου $480$ (βγαίνει $477{,}26$). Το $4\,773{,}3$ είναι δέκα φορές μεγαλύτερο.
\sol{16} $25$ φορές, γιατί $2{,}5:0{,}1=25$.
\sol{17} Η οθόνη δείχνει $0{,}14285714$ (οκτώ ψηφία). Στα εκατοστά: $0{,}14$. Τα ψηφία $142857$ επαναλαμβάνονται συνεχώς, επομένως $\frac{1}{7}=0{,}142857142857\ldots$.
\sol{18} Ελεύθερη απάντηση. Προσοχή να στοιχίσεις σωστά τα δεκαδικά και να ελέγξεις με εκτίμηση.
\sol{19} Ελεύθερη απάντηση. Το άθροισμα των θερμοκρασιών διαιρείται με το πλήθος τους, π.χ. $\frac{18{,}5+19+21+20{,}5+17+16{,}5+22}{7}=\frac{134{,}5}{7}\approx19{,}2$.
